\chapter{TD Control: SARSA, Q-Learning, and Exploration}

\textbf{Control} asks which \textit{action} to prefer from the current state, but one-step TD (5.2) only learns how good each \textit{state} is under a fixed policy---a table of $V$ does not score actions. With a known model, DP can still choose the action
\[
\pi(s)=\arg\max_a\sum_{s',r}P(s',r\mid s,a)\bigl[r+\gamma V(s')\bigr];
\]
without that $P$, the sum cannot be computed, so the $V$ table alone cannot rank actions. Model-free control therefore uses action-values $Q(s,a)$, which score each pair, so $\arg\max_a Q(s,a)$ needs no model. Monte Carlo control uses that idea, but learns $Q$ only from full returns $G_t$. This section keeps $Q$ for control and replaces that full-return backup by one-step TD, stepping $Q(S_t,A_t)$ toward a bootstrap from $Q$ at the next state. Two standard algorithms do that---\textbf{SARSA}, which is on-policy (the backup uses the next action the behavior policy actually takes), and \textbf{Q-learning}, which is off-policy (the backup uses a $\max$ over next actions).

\section{SARSA}

The update (5.2) steps $V(S_t)$ toward $R_{t+1}+\gamma V(S_{t+1})$, which cannot choose $\pi(S_t)$ without $P$. Replacing $V(S_t)$ by $Q(S_t,A_t)$ scores the action that was taken, by stepping $Q(S_t,A_t)$ toward $R_{t+1}+\gamma Q(S_{t+1},A_{t+1})$. Drawing both $A_t$ and $A_{t+1}$ from the same behavior $\pi$ makes the table estimate $Q^\pi$---that is \textbf{on-policy}. The method is \textbf{SARSA} (\textbf{S}tate--\textbf{A}ction--\textbf{R}eward--\textbf{S}tate--\textbf{A}ction), named for the five-tuple $(S_t,A_t,R_{t+1},S_{t+1},A_{t+1})$ that the update uses. SARSA was introduced under that name by Sutton (\citetxt{suttonbarto2018}{Sutton and Barto, 2018}); the learning update is a variation of the on-policy backup of \citetxt{rummery1994}{Rummery and Niranjan, 1994}, who studied it with neural function approximation. In this book we use the tabular form: one $Q(s,a)$ entry per pair, updated by (6.3) below.

The values $Q^\pi$ obey the compact Bellman expectation equation, the action-value analogue of (5.1) (\eqref{eq:B.18} in \appref{B}):

\begin{equation}
Q^\pi(s,a)
=
\mathbb{E}_\pi\bigl[R_{t+1}+\gamma Q^\pi(S_{t+1},A_{t+1})\mid S_t=s,A_t=a\bigr].
\tag{6.1}
\end{equation}

Here $A_{t+1}$ is drawn from $\pi$ at $S_{t+1}$. Observe one transition and that next action and, because $Q^\pi$ is unknown, put the current table estimate $Q$ in its place to form the one-step backup

\[
U_t=R_{t+1}+\gamma Q(S_{t+1},A_{t+1}).
\]

(If $S_{t+1}$ is terminal, take the $Q$ term as $0$.) The difference between that backup and the current estimate is the \textbf{TD error} on the pair---the same idea as (5.3):

\begin{equation}
\delta_t
=
\bigl(R_{t+1}+\gamma Q(S_{t+1},A_{t+1})\bigr) - Q(S_t,A_t).
\tag{6.2}
\end{equation}

Stepping $Q(S_t,A_t)$ toward $U_t$ with step-size $\alpha$ is (5.2) on pairs; the expanded and $\delta$ forms are the same update:

\begin{equation}
\begin{aligned}
Q(S_t,A_t)
&\leftarrow
Q(S_t,A_t)
+
\alpha\bigl[R_{t+1}+\gamma Q(S_{t+1},A_{t+1})-Q(S_t,A_t)\bigr]
\\
&=
Q(S_t,A_t)+\alpha\,\delta_t.
\end{aligned}
\tag{6.3}
\end{equation}

Because $A_{t+1}$ is chosen the same way as $A_t$, both from the behavior $\pi$, the actions in the backup match how we act, and the table estimates $Q^\pi$ for that same $\pi$. That is what \textbf{on-policy} means here. Recall that a \textbf{policy} is only the rule for picking an action given the state (Section 1.1), not the $Q$ table itself. Here that rule is $\varepsilon$-greedy in the \textbf{current} $Q$. With probability $1-\varepsilon$ take $\arg\max_a Q(s,a)$, otherwise pick uniformly in $\mathcal{A}$. Step 1 sets $Q(s,a)$ arbitrarily; $S_{t+1}$ comes from the environment after we take $A_t$; then $A_{t+1}$ is sampled from that $\varepsilon$-greedy $\pi$. As $Q$ improves, the same $\varepsilon$-greedy rule (now reading the new $Q$) is the updated behavior. That is evaluate--improve as in Monte Carlo control, but with a one-step TD evaluate step. The algorithm also returns $\pi'(s)=\arg\max_a Q(s,a)$. That $\pi'$ is greedy in the finished $Q$, not the $\varepsilon$-greedy $\pi$ used while learning.

Unlike Monte Carlo, SARSA does not generate the full episode first. Step 2 still loops for each episode, but the whole trajectory is not produced before updating. After $S_0$ and $A_0$, each inner step observes $(R_{t+1},S_{t+1})$, chooses $A_{t+1}$, and applies (6.3). The rest of the episode is generated only after that update, one transition at a time, until termination.

\begin{implnote}[Algorithm (SARSA: on-policy TD control)]
Input: $\varepsilon\in(0,1]$ (exploration); action set $\mathcal{A}$.

Output: $Q\approx Q^\pi$ for the $\varepsilon$-greedy $\pi$ used while learning; greedy $\pi'(s)=\arg\max_a Q(s,a)$.

\begin{enumerate}
\item Initialize $Q(s,a)$ arbitrarily, with $Q(\text{terminal},\,\cdot\,)=0$.
\item Loop for each episode:
\begin{itemize}
\item Initialize $S_0$ (starting state of the episode).
\item Choose $A_0$ from $Q$ $\varepsilon$-greedy in $S_0$ (over $\mathcal{A}$).
\item For $t=0,1,2,\ldots$ until termination:
\begin{itemize}
\item Take $A_t$, observe $R_{t+1}$, $S_{t+1}$.
\item Choose $A_{t+1}$ from $Q$ $\varepsilon$-greedy in $S_{t+1}$.
\item $Q(S_t,A_t)\leftarrow Q(S_t,A_t)+\alpha\bigl[R_{t+1}+\gamma Q(S_{t+1},A_{t+1})-Q(S_t,A_t)\bigr]$.
\end{itemize}
\end{itemize}
\item After enough episodes, return $Q$ and $\pi'(s)=\arg\max_a Q(s,a)$.
\end{enumerate}

\end{implnote}

\lstinputlisting{../code/ch06/sarsa.py}


\section{Q-learning}

The SARSA algorithm above returns $\pi'(s)=\arg\max_a Q(s,a)$. That $\pi'$ is greedy in the finished $Q\approx Q^\pi$; it is not the $\varepsilon$-greedy $\pi$ used while learning. $\pi^*$ is an optimal policy, $\pi^*(s)=\arg\max_a Q^*(s,a)$ (2.11). SARSA's $\pi'$ is not $\pi^*$ unless the exploring $\pi$ was already optimal. A different goal is to learn $Q^*$ so that this greedy $\pi'$ aims at $\pi^*$ while still exploring with behavior. That method is \textbf{Q-learning} (\citetxt{watkins1992}{Watkins and Dayan, 1992}), the same one-step TD pattern on pairs, except the bootstrap uses a $\max$ over next actions instead of the behavior's next action. The update (6.6) below is their tabular rule.

The targets $Q^*$ obey the compact Bellman optimality equation (\eqref{eq:C.5} in \appref{C}):

\begin{equation}
Q^*(s,a)
=
\mathbb{E}\bigl[R_{t+1}+\gamma\max_{a'} Q^*(S_{t+1},a')\mid S_t=s,A_t=a\bigr].
\tag{6.4}
\end{equation}

Unlike (6.1), the continuation after $S_{t+1}$ is a $\max$ over actions, not an action drawn from a behavior $\pi$. Observe one transition and, because $Q^*$ is unknown, put the current table estimate $Q$ in its place to form the one-step backup

\[
U_t=R_{t+1}+\gamma\max_{a'} Q(S_{t+1},a').
\]

(If $S_{t+1}$ is terminal, take the $\max$ term as $0$.) The \textbf{TD error} on the pair is again backup minus estimate:

\begin{equation}
\delta_t
=
\bigl(R_{t+1}+\gamma\max_{a'} Q(S_{t+1},a')\bigr) - Q(S_t,A_t).
\tag{6.5}
\end{equation}

Stepping $Q(S_t,A_t)$ toward $U_t$ is the same $\alpha$-update as (6.3), with this backup:

\begin{equation}
\begin{aligned}
Q(S_t,A_t)
&\leftarrow
Q(S_t,A_t)
+
\alpha\bigl[R_{t+1}+\gamma\max_{a'} Q(S_{t+1},a')-Q(S_t,A_t)\bigr]
\\
&=
Q(S_t,A_t)+\alpha\,\delta_t.
\end{aligned}
\tag{6.6}
\end{equation}

Unlike SARSA, where both the current action $A_t$ and the next action $A_{t+1}$ are chosen by the same behavior policy (e.g. $\varepsilon$-greedy in $Q$), Q-learning does not put a behavior-sampled $A_{t+1}$ in the backup. The next action is whichever $a'$ currently maximizes $Q(S_{t+1},a')$. Acting may still use $\varepsilon$-greedy for $A_t$ to explore; only how the backup treats the next action differs. That is why Q-learning is \textbf{off-policy}. No importance-sampling weights are needed. The $\max$ already puts the greedy target inside the one-step backup.

As in SARSA, the full episode is not generated first. Start at $S_0$, then for each step choose behavior $A_t$, observe $(R_{t+1},S_{t+1})$, and apply (6.6). There is no stored full return.

\begin{implnote}[Algorithm (Q-learning: off-policy TD control)]
Input: $\varepsilon\in(0,1]$ (behavior exploration); action set $\mathcal{A}$.

Output: $Q\approx Q^*$; greedy $\pi(s)=\arg\max_a Q(s,a)$.

\begin{enumerate}
\item Initialize $Q(s,a)$ arbitrarily, with $Q(\text{terminal},\,\cdot\,)=0$.
\item Loop for each episode:
\begin{itemize}
\item Initialize $S_0$ (starting state of the episode).
\item For $t=0,1,2,\ldots$ until termination:
\begin{itemize}
\item Choose $A_t$ from $Q$ $\varepsilon$-greedy in $S_t$ (behavior; over $\mathcal{A}$).
\item Take $A_t$, observe $R_{t+1}$, $S_{t+1}$.
\item $Q(S_t,A_t)\leftarrow Q(S_t,A_t)+\alpha\bigl[R_{t+1}+\gamma\max_{a'} Q(S_{t+1},a')-Q(S_t,A_t)\bigr]$.
\end{itemize}
\end{itemize}
\item After enough episodes, return $Q$ and $\pi(s)=\arg\max_a Q(s,a)$.
\end{enumerate}

\end{implnote}

\lstinputlisting{../code/ch06/q_learning.py}


\section{$\varepsilon$-greedy}

SARSA and Q-learning choose $A_t$ $\varepsilon$-greedy in the current $Q$ (SARSA also chooses $A_{t+1}$ that way). That rule is formulated here in mathematical form as a distribution $\pi(a\mid s)$ over $\mathcal{A}$. It is essentially the same choice as above, only written as (6.7), with an optional schedule (6.8) that shrinks $\varepsilon$. \textbf{Exploit} takes $\arg\max_a Q(s,a)$, so the action uses the current table. \textbf{Explore} draws uniformly from $\mathcal{A}$, so other pairs $(s,a)$ are visited and their $Q$ can update. $Q$ is updated only at the pair that was taken, so both are needed. With probability $1-\varepsilon$ exploit, and with probability $\varepsilon$ explore.

At state $s$, write $A^*=\arg\max_a Q(s,a)$ (break ties by any fixed rule). Because $A^*$ is still eligible on the uniform draw, its total probability is the greedy branch plus one share of that draw,

\[
\pi(A^*\mid s)
=
(1-\varepsilon)
+
\frac{\varepsilon}{\lvert\mathcal{A}\rvert}.
\]

Every other action appears only on the uniform draw, each with probability $\varepsilon/\lvert\mathcal{A}\rvert$. For $\varepsilon\in(0,1]$,

\begin{equation}
\pi(a\mid s)
=
\begin{cases}
1-\varepsilon+\dfrac{\varepsilon}{\lvert\mathcal{A}\rvert},
&
\text{if }a=A^*,
\\[0.75em]
\dfrac{\varepsilon}{\lvert\mathcal{A}\rvert},
&
\text{if }a\neq A^*.
\end{cases}
\tag{6.7}
\end{equation}

When $\varepsilon=0$, (6.7) is greedy. When $\varepsilon=1$, it is uniform. A fixed $\varepsilon>0$ explores forever, which helps coverage but keeps the behavior away from the greedy $\pi^*$. A common remedy is to shrink $\varepsilon$ over episodes $k=0,1,2,\ldots$ by a geometric factor. Writing $\varepsilon_k=\varepsilon_0^k$ does not work when the start is $\varepsilon_0=1$ (then $1^k=1$ forever), and it ties the starting level to the decay speed. Keep them separate: start at $\varepsilon_0\in(0,1]$ and multiply by $\lambda^k$ with $\lambda\in(0,1)$,

\begin{equation}
\varepsilon_k
=
\varepsilon_0\,\lambda^k.
\tag{6.8}
\end{equation}

Early episodes explore more; later ones concentrate on $\arg\max_a Q(s,a)$. In practice one may clip $\varepsilon_k$ below by a small floor so a little exploration remains; that floor is optional and is not part of (6.8). The same schedule can feed any method that chooses $A_t$ by (6.7)---including on-policy MC control, SARSA, and Q-learning.

The algorithm and code below are only that choice piece: given $Q$, $s$, and an $\varepsilon$ (optionally from (6.8)), return one action from (6.7). They are not a full control loop---no environment step, no $Q$ update. Plug \texttt{eps\_greedy} in wherever Chapters 4--6 chose $A_t$ (or SARSA's $A_{t+1}$) $\varepsilon$-greedy; use \texttt{eps\_schedule} when $\varepsilon$ should shrink over episodes.

\begin{implnote}[Algorithm (sample $A\sim\varepsilon$-greedy in $Q$)]
Input: $Q$, state $s$, $\varepsilon\in[0,1]$, action set $\mathcal{A}$.

\begin{enumerate}
\item $A^*\leftarrow\arg\max_a Q(s,a)$.
\item For each $a\in\mathcal{A}$, set $\pi(a\mid s)$ by (6.7): $1-\varepsilon+\varepsilon/\lvert\mathcal{A}\rvert$ if $a=A^*$, else $\varepsilon/\lvert\mathcal{A}\rvert$.
\item Return a draw from $\pi(\cdot\mid s)$.
\end{enumerate}

\end{implnote}

\lstinputlisting{../code/ch06/eps_greedy.py}


\section{Softmax}

$\varepsilon$-greedy in (6.7) uses $Q$ only to pick the greedy action $A^*$. Every other action then shares the same mass $\varepsilon/\lvert\mathcal{A}\rvert$, so a near-best action and a terrible one are treated alike when exploring.

To rank every action by $Q$ instead, use a \textbf{Boltzmann (Gibbs)} distribution over $\mathcal{A}$: $\pi(a\mid s)\propto\exp\bigl(Q(s,a)/\tau\bigr)$, so better actions are more likely and temperature $\tau>0$ sets how sharp the ranking is. \textbf{Softmax} writes that distribution in normalized form---exponentiate each score $Q/\tau$, then divide by the sum:

\begin{equation}
\pi(a\mid s)
=
\frac{\exp\bigl(Q(s,a)/\tau\bigr)}
{\sum_{b\in\mathcal{A}}\exp\bigl(Q(s,b)/\tau\bigr)}.
\tag{6.9}
\end{equation}

Large $\tau$ flattens $\pi$ toward uniform; small $\tau$ concentrates on $\arg\max_a Q(s,a)$. At $\tau=1$, (6.9) is the plain softmax of $Q$. Leave $\tau$ untuned and the sharpness tracks whatever scale $Q$ happens to have, so treat $\tau$ as a hyperparameter---fixed by trying a few values, or annealed over episodes like (6.8).

The algorithm and code below only return one draw from (6.9); they are not a full control loop. Use them in place of $\varepsilon$-greedy when selecting $A_t$ (or SARSA's $A_{t+1}$).

\begin{implnote}[Algorithm (sample $A\sim$ softmax in $Q$)]
Input: $Q$, state $s$, temperature $\tau>0$, action set $\mathcal{A}$.

\begin{enumerate}
\item For each $a\in\mathcal{A}$, set $u(a)\leftarrow\exp\bigl(Q(s,a)/\tau\bigr)$.
\item For each $a\in\mathcal{A}$, set $\pi(a\mid s)\leftarrow u(a)/\sum_{b\in\mathcal{A}} u(b)$ as in (6.9).
\item Return a draw from $\pi(\cdot\mid s)$.
\end{enumerate}

\end{implnote}

\lstinputlisting{../code/ch06/softmax.py}


SARSA, Q-learning, (6.7), and (6.9) all use a table of $Q(s,a)$: one entry per pair, updated only when that pair is taken. That table works because $\mathcal{S}$ so far is a small fixed set. The same $s$ returns, and the same cell is updated. When $\mathcal{S}$ is large or continuous, it runs into two problems: the table does not fit, and a later observed state is almost never that same $s$, so the stored number is not used again. Chapter 7 replaces the table by a function $\hat{Q}(s,a;\mathbf{w})$ of a weight vector $\mathbf{w}$. Only those weights are stored and updated. SARSA and Q-learning keep the same backups, and (6.7) and (6.9) still choose $A_t$, now from $\hat{Q}$.
